\subsection{Компьютерное зрение: ArUco и AprilTag}

Для роботов существуют специальные визуальные метки (маркеры), которые легко распознать на изображении. Самые популярные — \textbf{ArUco} и \textbf{AprilTag}.

\begin{warnbox}[Начиная с этой главы нужен настоящий дрон]
Все последующие главы курса рассчитаны на реальный Pioneer с камерой и распечатанными маркерами. Встроенный в браузер симулятор дает настоящее видео с виртуальной камеры, но не реализует компьютерное зрение: \texttt{detectMarkers()} там всегда возвращает пустой список (маркер не будет найден никогда), \texttt{cv2.solvePnP()} всегда сообщает о неудаче, а библиотеки \texttt{mediapipe} в браузере нет вообще --- \texttt{import mediapipe} завершится ошибкой \texttt{ModuleNotFoundError}. Поэтому код из этой и следующих глав синтаксически верен и пригоден для реального дрона, но запуск его в браузерном редакторе Python даст либо <<тишину>> (маркеры не находятся), либо ошибку импорта. Это ограничение среды, а не ошибка в вашей программе.
\end{warnbox}

\subsubsection*{Что такое ArUco-маркер?}
Это черно-белый квадрат с уникальным узором внутри. У каждого маркера есть свой \textbf{ID} (номер).
\begin{itemize}
    \item \textbf{Словарь (Dictionary):} Набор допустимых маркеров. Pioneer обычно использует словарь \texttt{DICT\_6X6\_50} (размер сетки 6x6 бит, всего 50 уникальных маркеров).
    \item \textbf{AprilTag:} Другой популярный формат, поддерживаемый той же библиотекой (например, \texttt{DICT\_APRILTAG\_36h11}).
\end{itemize}

\subsubsection*{Детекция маркеров}
Для работы с камерой и маркерами мы используем библиотеку \textbf{OpenCV} (\texttt{cv2}).

\begin{codefile}[python]{python/examples/aruco_examples/11_aruco_basics_example1_detect.py}
\end{codefile}

\subsubsection*{Разбор кода}
1. \texttt{camera.get\_cv\_frame()} получает изображение с камеры дрона.
2. \texttt{cv2.cvtColor} переводит картинку в ч/б формат (так алгоритм работает быстрее и надежнее).
3. \texttt{detectMarkers} ищет квадраты и сверяет их узор со словарем.
4. Возвращает:
   \begin{itemize}
       \item \texttt{corners}: Координаты углов найденных маркеров.
       \item \texttt{ids}: Список номеров (ID) найденных маркеров.
   \end{itemize}

\begin{taskbox}[Охота на маркеры]
Распечатайте несколько маркеров ArUco (или откройте их на экране телефона). Напишите программу, которая пишет в консоль <<Вижу маркер X!>>, где X — номер маркера.
\end{taskbox}

\subsubsection*{Определение координат маркера}
Зная размер маркера и параметры камеры (полученные при калибровке), можно вычислить расстояние и ориентацию относительно маркера.

\begin{codefiletitle}[python]{Определение координат}{python/examples/aruco_examples/11_aruco_basics_example2_coordinates.py}
\end{codefiletitle}

\subsubsection*{Удержание позиции по маркеру}
В этом примере дрон использует координаты маркера для удержания позиции над ним.

\begin{codefiletitle}[python]{Полет по маркеру}{python/examples/aruco_examples/11_aruco_basics_example3_flight.py}
\end{codefiletitle}
